\chapter{Concluzii}
\label{cap:concluzii}

\section{Concluzii generale}

Lucrarea a pornit de la o întrebare pragmatică — se poate construi, din componente accesibile, un robot capabil să susțină activități terapeutice reale pentru copii neurodivergenți? — iar răspunsul demonstrat practic este afirmativ. TriSense există, funcționează cap-coadă pe hardware real și acoperă un ciclu complet de interacțiune: copilul apasă un buton pe robot, vorbește, robotul îl înțelege, îi răspunde cu voce și mișcare, îi propune jocuri, îi verifică vizual construcțiile LEGO și consemnează totul pentru terapeut.

Din perspectivă tehnică, principala concluzie este că arhitectura distribuită pe trei niveluri — fiecare dispozitiv făcând strict ceea ce știe cel mai bine — a fost decizia care a făcut restul posibil. Hub-ul LEGO nu știe că există internet; ESP32 nu știe ce spune modelul de limbaj; PC-ul nu atinge niciun motor. Această separare a făcut fiecare componentă testabilă izolat și a localizat imediat fiecare defect — de la redarea sacadată (rezolvată prin buffering pe ESP32) până la interferența modelului generativ cu logica jocului (rezolvată prin izolarea stărilor în creier).

Din perspectivă aplicativă, sistemul transpune fidel principiile documentate în literatura de specialitate: predictibilitatea și repetabilitatea pe care studiile despre roboți și autism \cite{scassellati2012, dautenhahn2004} le identifică drept cheia disponibilității copilului; materialul LEGO, intrinsec motivant, validat de cercetările LBT \cite{legoff2004, owens2008}; feedback-ul imediat și exclusiv pozitiv; și rolul de mediator, nu de substitut — robotul susține sesiunea, terapeutul o conduce.

Merită consemnată și o concluzie metodologică despre folosirea modelelor de inteligență artificială generativă într-un sistem cu reguli stricte: puterea lor (flexibilitate, limbaj natural, verificare vizuală fără antrenare) vine la pachet cu nevoia de a le îngrădi explicit — prin prompturi atent formulate, prin praguri calibrate empiric și, mai ales, prin protejarea fluxului de control al jocurilor de creativitatea modelului conversațional.

\section{Direcții de dezvoltare ulterioară}

Catalogul de funcționalități schițat în faza de proiectare conține mult mai multe procedee decât au încăput în implementarea curentă, iar câteva direcții se conturează natural:

\begin{itemize}
    \item \textbf{validarea clinică extinsă}: primele sesiuni reale cu un copil neurodivergent (capitolul~\ref{cap:evaluare}) au decurs încurajator, copilul interacționând firesc cu robotul pe parcursul activităților; pasul rămas este confirmarea eficacității printr-un studiu amplu — eșantion de copii, protocol formal și avize etice — în colaborare cu terapeuți specializați în LBT, fiindcă abia astfel de date pot transforma argumentele de plauzibilitate din această lucrare în dovezi;
    \item \textbf{jocurile aflate încă în stadiu de prototip}: N-back verbal (actualizarea memoriei de lucru) și Stroop adaptat (inhibiția interferenței), ambele deja proiectate pe hardware-ul existent — completarea lor ar acoperi întreg tripticul funcțiilor executive descris de Diamond;
    \item \textbf{detecția stării emoționale din voce}: analiza prozodiei (ton, ritm, intensitate) ar permite robotului să sesizeze frustrarea înainte ca ea să escaladeze și să propună proactiv pauza de respirație;
    \item \textbf{înlocuirea ferestrei fixe de înregistrare} cu detecția automată a activității vocale (VAD), pentru un dialog mai natural;
    \item \textbf{funcționarea offline}: modele locale de transcriere și dialog ar elimina dependența de internet — relevantă pentru cabinete cu conectivitate slabă;
    \item \textbf{interfața terapeutului}: un panou web pentru configurarea sesiunii (jocuri, niveluri, durată), urmărirea în timp real și vizualizarea istoricului de metrici per copil;
    \item \textbf{personalizarea pe termen lung}: memoria persistentă a sistemului (preferințe, progres, nivel atins) poate alimenta adaptarea automată a dificultății între sesiuni.
\end{itemize}

Închei cu observația care mi se pare cea mai valoroasă din întregul proiect: niciuna dintre componentele folosite — setul LEGO, plăcile ESP32, serviciile de inteligență artificială — nu este exotică sau scumpă. Tot ce desparte un set educațional de un instrument terapeutic este integrarea. Iar integrarea, spre deosebire de hardware-ul specializat, se poate replica oriunde există un calculator, o rețea Wi-Fi și cineva dispus să o facă.
